% 夜半用户手册 · 第 1 章样章（中文）
% T-DOC-001 第一批交付：大纲＋术语表＋第一章样章。
% 事实来源：README.md、docs/manual/01-overview.md、02-projects.md、
% crates/yeban-app/src/help.rs（引导步骤）、cli.rs 的 SHORTCUTS。

\chapter{快速上手}
\label{ch:quickstart}

夜半（Yeban）是一款开源数字音频工作站：纯 Rust 实时音频引擎之上是一套 Slint
原生图形界面，并对外提供完整的 MCP 工具面。人和 AI 可以在同一个工程里并肩工作。
本章带你从安装走到第一次导出，全程只用它真实存在的功能。

\section{安装与首次启动}

\subsection{选择安装方式}

最省事的是仓库 \texttt{README.md} 的 Download 表所列的 nightly 安装包（跟随
\texttt{main} 分支，未签名）：Windows x86\_64 用 MSI，macOS 用 DMG（Apple Silicon
与 Intel 各一份），Linux 用 DEB（x86\_64 与 arm64 各一份）。便携版
\texttt{.tar.gz} 与 \texttt{SHA256SUMS} 在 nightly 发布页。

\paragraph{从源码构建} 需要 Rust 1.99.0（由 \texttt{rust-toolchain.toml} 固定，
\texttt{rustup} 会自动安装）、C++ 编译器与 \texttt{libclang}
（\texttt{yeban-dsp} 用 bindgen 生成绑定）。Ubuntu/Debian 上的依赖与构建命令：

\begin{verbatim}
sudo apt-get install -y clang libclang-dev pkg-config libfontconfig1-dev \
  libfreetype-dev libxkbcommon-dev libxkbcommon-x11-0 libwayland-dev \
  libx11-dev libgl1-mesa-dev libasound2-dev
cargo build --release -p yeban-app
\end{verbatim}

运行用 \texttt{cargo run --release -p yeban-app}；没有显示器的机器加
\texttt{-- --headless}。内存紧张时用
\texttt{cargo build --profile release-small -p yeban-app}。

\begin{figure}[htbp]
\centering
\begin{tikzpicture}[node distance=1.1cm, >=stealth,
    box/.style={draw, rounded corners, align=center, minimum width=2.2cm, minimum height=0.8cm}]
  \node[box] (a) {安装包 / 源码构建};
  \node[box, right of=a] (b) {首次启动};
  \node[box, right of=b] (c) {新手引导};
  \node[box, right of=c] (d) {建轨与放片段};
  \node[box, right of=d] (e) {播放试听};
  \node[box, right of=e] (f) {导出 WAV};
  \draw[->] (a) -- (b); \draw[->] (b) -- (c); \draw[->] (c) -- (d);
  \draw[->] (d) -- (e); \draw[->] (e) -- (f);
\end{tikzpicture}
\caption{第一章的动手路线：安装 $\to$ 启动 $\to$ 引导 $\to$ 编排 $\to$ 试听 $\to$ 导出}
\label{fig:ch01-flow}
\end{figure}

\subsection{首次启动与新手引导}

第一次启动时，应用会自动打开新手引导（落盘状态判定为「该弹」才弹；点过「不再
显示」或已走完一遍就不会再自动出现）。引导带你认界面：欢迎页、走带栏、轨道区、
钢琴卷帘、混音控制台，最后是检查器与音色库。每一步都会把界面切到正在讲的那一区。
之后想再看，从「帮助 $\to$ 新手引导」重开。

\screenshot{ch01-onboarding-tour}{首次启动自动弹出的新手引导卡片，停留在「走带栏」一步}

\section{界面总览}

\subsection{菜单栏与走带栏}

顶部菜单栏共八组：文件、编辑、轨道、走带、混音、视图、诊断、帮助。走带栏在
窗口最上方：播放 / 暂停（空格）、停止、归零到 \texttt{001.01.000}、循环开关
（C）、节拍器与预备拍，以及速度（BPM）显示与敲入。走带栏右侧是工作区切换：
全景、现场激发、线性编曲、混音控制台，数字键 1--4 对应。

\screenshot{ch01-main-window}{内置演示工程打开后的主窗口，标出菜单栏、走带栏、编曲区、检查器、底部控制台}

\subsection{中部：编曲区}

一条轨道一行：左边是轨道头（名称、静音 / 独奏、音量），右边时间轴上摆着片段。
片段可以拖动、拉边缘改长度、复制（Cmd/Ctrl+C / Cmd/Ctrl+V）、剪切
（Cmd/Ctrl+X）与删除。上方是小节标尺与段落标记；下方可展开和弦轨与自动化泳道。
按 Tab 打开 Clip 激发矩阵（Session 视图），按格子触发片段。

\subsection{左侧：检查器与音色库}

左栏上半是通道检查器，显示当前选中轨的通道条与设备参数，按 I 显隐；下半是
音色库，浏览预置音色与 SFZ 音色库。轨道头工具行最前面是「轨道菜单」按钮，
汇集本轨的操作。

\subsection{底部控制台与边缘}

底部三个页签：钢琴卷帘（P / F7）、混音控制台（X / F8）、设备机架（F9）；
F3 把控制台最大化 / 还原。右侧可打开监听控制室（F4 显隐），显示主控监听与
响度读数。最下方状态栏给出选区信息、实时和弦识别、快捷键提示与音频设备状态；
再往下是全曲总览面板，用一张图看全曲能量与结构。把鼠标停在控件上会出现悬浮
提示，按 T 整体开关。

\begin{table}[htbp]
\centering
\caption{第一章涉及的键位（规范写法；macOS 为 Cmd+，其余为 Ctrl+）}
\label{tab:ch01-keys}
\begin{tabular}{ll}
\hline
键位 & 作用 \\
\hline
Space & 播放 / 暂停 \\
Shift+Space & 从光标处续播 \\
C & 循环播放开关 \\
Enter & 回到起点 \\
Tab / F5 & 视图切换 / Session 视图 \\
F6 & Arrangement 视图 \\
P / F7 & 钢琴卷帘 \\
X / F8 & 混音控制台 \\
F9 & 设备机架 \\
I & 显隐左侧检查器 \\
Cmd/Ctrl+S & 保存工程 \\
Cmd/Ctrl+Shift+S & 另存为 \\
Cmd/Ctrl+Z / Cmd/Ctrl+Shift+Z & 撤销 / 重做 \\
Cmd/Ctrl+T & 在播放头处切分选中的剪辑 \\
\hline
\end{tabular}
\end{table}

\section{第一个工程}

\subsection{新建，或先用演示工程逛一圈}

工程文件的扩展名是 \texttt{.yeban}，它是唯一的工程格式：一个容器里装着轨道、
片段、音符、设备链与所用素材。

「文件 $\to$ 新建空白工程」得到一个只有主总线的空工程。如果启动时不指定工程，
应用会打开一个演示工程（鼓、贝斯、铺底、主音、弦乐、打击乐），方便先逛界面；
它不是你的作品文件，随便改不会弄丢任何东西。

\subsection{建轨、放片段、播放}

在空工程里，用轨道菜单新建一条音轨；打开音色库挑一个预置音色或 SFZ 音色，
把音频片段放到时间轴上。按空格播放，走带栏显示当前位置与速度；需要反复听一段
时按 C 打开循环并框选区间。片段放歪了就拖动改位置、拖边缘改长度；要在播放头
处切开，选中片段按 Cmd/Ctrl+T。

\screenshot{ch01-first-clip}{空工程里第一条轨道上放好一个片段，播放头停在片段中部}

\subsection{保存与导出}

「文件 $\to$ 保存工程（Cmd/Ctrl+S）」把当前工程写回原文件；
「文件 $\to$ 另存为（Cmd/Ctrl+Shift+S）」换名字或位置。保存是原子写入，
写一半断电也不会留下半个文件。

导出整轨 WAV 用「文件」菜单里的导出入口；分轨导出（Stems）会把每条源节点轨
弹成一个 32-bit Float WAV，写进固定输出目录，并给出一份 \texttt{stems-manifest.json}
索引与一行导出回执。主总线不出文件，但会在跳过清单里点名——这不是丢轨。

\screenshot{ch01-export-dialog}{导出对话框与导出后的回执/toast 提示}

\section{帮助主题映射}

\begin{table}[htbp]
\centering
\caption{本章小节 $\to$ 应用内帮助锚点}
\label{tab:ch01-help}
\begin{tabular}{ll}
\hline
本章小节 & 帮助主题 \\
\hline
安装与首次启动 & \texttt{getting-started-install} \\
首次启动与新手引导 & \texttt{getting-started-onboarding} \\
菜单栏与走带栏 & \texttt{interface-transport} \\
中部：编曲区 & \texttt{interface-arrangement} \\
左侧：检查器与音色库 & \texttt{interface-inspector-library} \\
底部控制台与边缘 & \texttt{interface-console} \\
新建，或先用演示工程 & \texttt{project-new-demo} \\
建轨、放片段、播放 & \texttt{project-first-arrange} \\
保存与导出 & \texttt{project-save-export} \\
\hline
\end{tabular}
\end{table}

\paragraph{本章学完能做什么} 独立走通「安装 $\to$ 启动 $\to$ 建轨 $\to$ 放片段 $\to$ 播放 $\to$ 导出」全流程，并能说出界面五个主要区域各管什么。

\paragraph{动手练习} 新建一个空工程，在轨道 1 放一段演示音频片段；按空格播放并用 C 循环两小节；用 Cmd/Ctrl+S 保存；导出整轨 WAV 与分轨各一次；最后用「文件 $\to$ 打开 / 恢复工程」重新打开，核对片段还在。
